
\begin{Exercice}[5 minutes]\textbf{Maxim}
En Java, quel code définit correctement un getter pour un attribut privé type?

\begin{lstlisting}[style=javastyle]
// Option A
public void getType() { return type; }

// Option B
private String getType() { return this.type; }

// Option C
public String getType() { return type; }

// Option D
public static String getType() { return type; }
\end{lstlisting}

\begin{solution}
\texttt{public String getType() \{ return type; \}}
\end{solution}

\end{Exercice}

\begin{Exercice}[5 minutes]\textbf{Maxim}
En Java, quelle est la syntaxe correcte pour indiquer qu’une classe \texttt{VoitureElectrique} hérite de \texttt{Voiture}?

\mcqlist{
    \texttt{class VoitureElectrique -> Voiture},
    \texttt{class VoitureElectrique extends Voiture},
    \texttt{class VoitureElectrique inherits Voiture},
    \texttt{class VoitureElectrique super Voiture}
}
\begin{solution}
\texttt{class VoitureElectrique extends Voiture}
\end{solution}

\end{Exercice}

\begin{Exercice}[5 minutes]\textbf{Léa}
Soit la classe suivante:

\begin{lstlisting}[style=javastyle]
public class Cat {
    public String name;
    private int age;
    private int mood = 5;

    public Cat(String name, int age) {
        this.name = name;
        this.age = age;
    }

    public String getName() {
        return this.name;
    }
    public int getAge() {
        return age;
    }
    public int getMood() {
        return mood;
    }
    public void setAge(int age) {
        this.age = age;
    }
    public void setName(String name) {
        this.name = name;
    }
    public void setMood(int mood){
        this.mood = mood;
    }

    public boolean canPlay(Cat otherCat){
        ???
    }
}
\end{lstlisting}

La méthode \texttt{canPlay} décrit si 2 chats peuvent jouer ensemble. Deux chats ne veulent pas jouer ensemble si le mood de l'un des deux est inférieur (strictement) à 3, dans ce cas la fonction retourne faux, sinon elle retourne vrai.\newline 

Ecrivez le code de la fonction \texttt{canPlay}.

\begin{solution}
\lstinputlisting[style=javastyle]{code/Q66.java}
\end{solution}
\end{Exercice}

\begin{Exercice}[5 minutes]\textbf{Léa}
Soit les 2 classes suivantes:

\begin{lstlisting}[style=javastyle]
public class Animal {
    private String type;
    protected String name;
    protected int age;
    protected int nb_animal = 0;

    public Animal(String type, String name, int age){
        this.type = type;
        this.name = name;
        this.age = age;
    }
}
\end{lstlisting}

\begin{lstlisting}[style=javastyle]
public class Cat extends Animal {
    #Ligne 1 à remplacer

    public Cat(String name, int age, String type, int mood) {
        #Ligne 2
        #Ligne 3
    }

    public String getName() {
        return this.name;
    }
    public int getAge() {
        return age;
    }
    public int getMood() {
        return mood;
    }
    public void setAge(int age) {
        this.age = age;
    }
    public void setName(String name) {
        this.name = name;
    }
    public void setMood(int mood){
        this.mood = mood;
    }
}
\end{lstlisting}

Compléter le code suivant en ajoutant l'attribut nécessaire à la ligne 1 et compléter le constructeur à la ligne 2 et 3.

\begin{solution}
Ligne 1: \texttt{private int mood;} \newline 
Ligne 2: \texttt{super(type,name,age);} \newline 
Ligne 3: \texttt{this.mood = mood;}
\end{solution}


Soit la suite donnée par $u_0 = 4 \quad \text{et} \quad \forall n > 1,\; u_n = n + u_{n-1}^2$.

\begin{lstlisting}[style=pythonstyle]
def suite(n):
  if n== 0:
    return ???
  else:
    return ???
\end{lstlisting}

Par quoi faut-il remplacer les ??? pour que le code correspond à l'expression donnée plus-haut?

\mcqlist{
    {4, $n + suite(n+1)$},
    {0, $n + u(n-1)**2$},
    {4, $n + suite(n-1)**2$},
    {0, $n**2 + suite(n+1)$}
}

\begin{solution}
4, $n + suite(n-1)**2$
\end{solution}

\end{Exercice}


\begin{Exercice}[5 minutes]\textbf{Sean}

% Quelle sera la sortie dans la console après l'exécution du programme \texttt{CoffeShop}?

% \begin{lstlisting}[style=javastyle]
% // Fichier: CoffeeMachine.java
% public class CoffeeMachine {
%   public static int coffeesMadeCount = 0;
%   public CoffeeMachine() {}
%   public void makeCoffee() {
%       CoffeeMachine.coffeesMadeCount++;
%   }
% }
% \end{lstlisting}

% \begin{lstlisting}[style=javastyle]
% // Fichier: CoffeeShop.java
% public class CoffeeShop {
%   public static void main(String[] args) {
%     CoffeeMachine machine1 = new CoffeeMachine();
%     CoffeeMachine machine2 = new CoffeeMachine();
%     machine1.makeCoffee();
%     System.out.println(machine1.coffeesMadeCount);
%     machine2.makeCoffee();
%     System.out.println(machine2.coffeesMadeCount);
%     CoffeeMachine.coffeesMadeCount = 0;
%     System.out.println(machine1.coffeesMadeCount);
%     }
% }
% \end{lstlisting}

% \mcqlist{
%      1 \; 1 \; 0,
%      1 \; 2 \; 0,
%      1 \; 1 \; 1,
%      1 \; 2 \; 1
% }


% \begin{solution}
% 1 \; 2 \; 0
% \end{solution}

% Comment empêcher le CoffeeShop de réinitialiser l'attribut CoffeesMadeCount, tout en laissant le CoffeeShop consulter la valeur de l'attribut, sélectionnez toutes les modifications nécessaires?

% \mcqlist{
%     Rendre CoffeesMadeCount privée,
%     Créer un getter pour CoffeesMadeCount,
%     Créer un setter pour CoffeesMadeCount,
%     Supprimer le mot-clé static de CoffeesMadeCount
% }

% \begin{solution}
% Rendre CoffeesMadeCount privée et créer un getter pour CoffeesMadeCount.
% \end{solution}

Parmi les affirmations suivantes, lesquelles sont corrects: 

\mcqlist{
{En Java, les attributs d'une classe peuvent être initialisés directement lors de leur déclaration, en dehors de tout constructeur.},
{En Python, des attributs propres à chaque instance peuvent être déclarer en dehors de toute méthode (dans le corps de la classe)},
{
    En Python, la référence à l'instance (self) doit être passée explicitement comme premier argument de chaque méthode d'instance, contrairement au mot-clé this en Java qui est implicite.
},
{
    En Python, le constructeur d'une classe est défini par la méthode spéciale \_\_init\_\_, alors qu'en Java, il porte le même nom que la classe.
}
}

\begin{solution}
A, C, et D.
\end{solution}


\end{Exercice}

\begin{Exercice}[5 minutes]\textbf{Amina}
Qu'affiche le programme suivant?

\begin{lstlisting}[style=pythonstyle]
class Logger:
    def __init__(self):
        self.events = []
    def log(self, msg):
        self.events.append(msg)
    def summary(self):
        return "-".join(self.events)

class Process(Logger):
    def __init__(self, name):
        super().__init__()
        self.name = name
    def start(self):
        self.log("start:" + self.name)
        self.execute()
    def execute(self):
        self.log("exec:" + self.name)
    def stop(self):
        self.log("stop:" + self.name)

class TimedProcess(Process):
    def __init__(self, name, ticks):
        super().__init__(name)
        self.ticks = ticks
    def execute(self):
        for i in range(self.ticks):
            self.log(f"tick{i}")
        super().execute()
        self.stop()

p = TimedProcess("P1", 2)
p.start()
print(p.summary())
\end{lstlisting}

\mcqlist{
    start:P1-tick0-tick1-stop:P1,
    tick0-tick1-start:P1-exec:P1-stop:P1,
    start:P1-tick0-tick1-exec:P1-stop:P1,
    Ce programme n'affiche rien.
}
\begin{solution}
start:P1-tick0-tick1-exec:P1-stop:P1
\end{solution}

Qu'affiche le code suivant?

\begin{lstlisting}[style=javastyle]
class Enseignant {
    protected int baseSalary = 3000;
    public int salaireMensuel() { return baseSalary; }
}

class Professeur extends Enseignant {
    private int committees;
    public Professeur(int committees) { this.committees = committees; }
    public int salaireMensuel() { return super.salaireMensuel() + committees * 200; }
}

class Collaborateur extends Enseignant {
    private int hours, rate;
    public Collaborateur(int hours, int rate) { this.hours = hours; this.rate = rate; }
    public int salaireMensuel() { return hours * rate; }
}

public class Test {
    public static void main(String[] args) {
        Enseignant e = new Professeur(3);
        Enseignant f = new Collaborateur(10, 50);
        System.out.println(e.salaireMensuel() + f.salaireMensuel());
        e = f;
        System.out.println(e.salaireMensuel());
    }
}
\end{lstlisting}

\mcqlist{
    4100 \newline 
    500,
    3500 \newline 
    3000,
     4100 \newline 
    3000,
    3000 \newline 
    500
}

\begin{solution}
4100 \newline 
500
\end{solution}

Laquelle des affirmations suivantes est vraie?

\mcqlist{
    {L’héritage permet à une sous-classe d’accéder directement aux champs private de sa superclasse, ce qui favorise la modularisation.},
{L’encapsulation vise à cacher la complexité interne d’un objet, permettant de l’utiliser uniquement via les méthodes définies dans sa classe.},
 {La modularisation consiste à regrouper dans une seule classe toutes les fonctionnalités nécessaires au programme pour éviter la duplication de code.},
 {Le sous-typage empêche une instance d’une sous-classe d’être utilisée là où un objet de la classe mère est attendu}.
}

\begin{solution}
L’encapsulation vise à cacher la complexité interne d’un objet, permettant de l’utiliser uniquement via les méthodes définies dans sa classe.
\end{solution}

\end{Exercice}

\begin{Exercice}[5 minutes]\textbf{Léa}
Soit le code ci dessous, par quoi faut-il remplacer les ??? pour que le code fonctionne?

\begin{lstlisting}[style=pythonstyle]
class Circle:
    def __init__(self,rayon):
        self.__rayon = rayon

    def __str__(self):
        return f"Ce cercle a un rayon de {self.__rayon}"
    def get_rayon(self):
        return self.__rayon

    def __eq__(self, other):
        if ???:
            if self.__rayon == other.get_rayon():
                return True
            else: 
                return False
  else:
    return False
\end{lstlisting}

\mcqlist{
    \texttt{isinstance(other, Circle)},
    \texttt{isinstance(other,self)},
    \texttt{other.isCircle()},
    \texttt{other.equal(Circle)}
}

\begin{solution}
\texttt{isinstance(other, Circle)}
\end{solution}

Qu'affiche le code ci dessous?
\begin{lstlisting}[style=pythonstyle]
counter = 0
class Vehicle:
    def __init__(self, brand):
        global counter
        counter += 1
        self._id = counter
        self._brand = brand

    def start(self):
        print("Starting vehicle-" + str(self._id))
        self.engine()

    def engine(self):
        print("- generic engine started")

class Car(Vehicle):
    def __init__(self, brand, doors):
        super().__init__(brand)
        self.doors = doors

    def start(self):
        print("Starting car-" + str(self._id) +" " + self._brand)
        super().start()

    def engine(self):
        print("- car engine with " + str(self.doors) + " doors started")

vehicle = Car("Toyota", 4)
vehicle.start()

vehicle = Vehicle("Generic")
vehicle.start()
\end{lstlisting}


\mcqlist{
{Starting car-1 \newline 
Starting vehicle-1 \newline 
- car engine with 4 doors started \newline 
AttributeError: 'Vehicle' object has no attribute 'engine'},
{
 Starting car-1 Toyota \newline 
Starting vehicle-1 \newline 
- generic engine started \newline 
Starting vehicle-2 \newline 
- generic engine started \newline 
},
{
Starting car-1 Toyota  \newline 
Starting vehicle-1 \newline 
- car engine with 4 doors started \newline 
Starting vehicle-2 \newline 
- generic engine started
},
{
Starting car-1 \newline 
- car engine with 4 doors started \newline 
Starting vehicle-2 \newline 
- generic engine started
},
}
\begin{solution}
Option C: \newline 

Starting car-1 Toyota  \newline 
Starting vehicle-1 \newline 
- car engine with 4 doors started \newline 
Starting vehicle-2 \newline 
- generic engine started
\end{solution}


\end{Exercice}


\begin{Exercice}[5 minutes]\textbf{Aurore}
Laquelle des phrases suivantes est fausse concernant les attributs?

\mcqlist{
{Le mot clé static en Java désigne un attribut de classe},
{\texttt{self.\_\_attribut} désigne un attribut privé en Python},
{Un getter est nécessaire pour utiliser un attribut privé d'un objet en dehors de la classe où l'attribut est défini},
{Le mot clé self est nécessaire pour utiliser un attribut d'un objet sur lequel une méthode est appelée en Java}
}

\begin{solution}
Option D: Le mot clé self est nécessaire pour utiliser un attribut d'un objet sur lequel une méthode est appelée en Java.
\end{solution}

Les trois classes java suivantes sont définies:

\begin{lstlisting}[style=javastyle]
//classe 1
public class Vehicle {
    public String name;
    public int speed;
    public int nb_wheels;
    static int nb_vehicles;
    public Vehicle(String name,int speed,int nb_wheels){
        this.name=name;
        this.speed=speed;
        this.nb_wheels=nb_wheels;
        nb_vehicles++;
    }
}

//classe 2
public class Car extends Vehicle {
    public boolean isOpenRoof;
    public Car(String name, int speed,boolean or){
        super(name,speed, 4);
        this.isOpenRoof=or;
    }
}

//classe 3
public class Bike extends Vehicle {
    public float height;
    public Bike(String name, int speed,float height){
        super(name,speed, 2);
        this.height = height;
    }
}
\end{lstlisting}

Quelle ligne ne provoque aucune erreur?
\mcqlist{
  \texttt{Vehicle train = new Vehicle("train", 100, 30.6);},
  \texttt{Vehicle jeep = new Car("jeep",140,false);},
  \texttt{Bike kawa = new Bike("kawa", 300,true);},
  \texttt{Car lambo = new Vehicle("lambo",500,4);}
}


\begin{solution}
\texttt{Vehicle jeep = new Car("jeep",140,false);}
\end{solution}

\end{Exercice}

\begin{Exercice}[5 minutes]\textbf{Aurore}
Parmi les expressions suivantes, lesquelles retournent \texttt{true} lorsque les objets a et b ont les mêmes attributs et sous quelle condition?

\mcqlist{
  {\texttt{a.equals(b) (sans condition particulière)}},
  {\texttt{a.equals(b) (si la méthode equals() est redéfinie dans ce but)}},
  {\texttt{a == b (si a et b sont de type int)}},
  {\texttt{a == b (si on a assigné la valeur a à b ou l'inverse)}},
  {\texttt{aucune de ces réponses}}
}

\begin{solution}
B, C, D. 
\end{solution}

Soit la classe suivante en Java.

\begin{lstlisting}[style=javastyle]
public class Car extends Vehicle {
    private int  nb_seats;
    public boolean isOpenRoof;
    private static float nb_mini;
    public double height;
    public Car(String name, int speed, int nb_seats, boolean or, double height){
        super(name,speed);
        this.nb_seats = nb_seats;
        this.isOpenRoof=or;
        this.height = height;
        if (nb_seats<5){
            nb_mini++;
        }
    }
}
\end{lstlisting}
Lesquelles des affirmations suivantes sont vraies?
\mcqlist{
  {Les attributs name et speed peuvent être définis comme privés ou publics dans la classe mère},
  {La valeur par défaut de \texttt{or} est false},
  {Les attributs \texttt{nb\_seats} et \texttt{nb\_mini} sont privés},
  {Les attributs \texttt{nb\_mini} et \texttt{height} sont de même type},
  {L'attribut \texttt{nb\_mini} est un attribut de classe},
  {La classe mère défini exactement deux attributs}
}

\begin{solution}
C, E.
\end{solution}

\end{Exercice}

\begin{Exercice}[5 minutes]\textbf{Maxim}
Par quelle valeur faut-il remplacer \texttt{\#LINE REMOVED} afin que le code exécute correctement?

\begin{lstlisting}[style=pythonstyle]
class Person:
    def __init__(self, name):
        #LINE REMOVED
p = Person("Alice")
print(p.name)
\end{lstlisting}

\mcqlist{
    {\texttt{self.name = name }},
    {\texttt{name = self.name}},
    {\texttt{return name}},
    {\texttt{print(name)}}
}

\begin{solution}
Option A.
\end{solution}

Qu'affiche le programme suivant?

\begin{lstlisting}[style=pythonstyle]
class Compteur:
    def __init__(self, valeur):
        self.valeur = valeur

    def ajouter(self, x):
        self.valeur += x

class Multiplicateur:
    def __init__(self, facteur):
        self.facteur = facteur

    def multiplier(self, x):
        return x * self.facteur

c = Compteur(5)
m = Multiplicateur(3)

c.ajouter(2)
resultat = m.multiplier(c.valeur)

print(resultat)
\end{lstlisting}

\mcqlist{15, 21, 17, 7}
\begin{solution}
21
\end{solution}

\end{Exercice}


\begin{Exercice}[5 minutes]\textbf{Sean}


Soit cet extrait de code qui compile et s'exécute sans erreurs. Lesquelles des affirmations suivantes pourraient décrire avec précision la relation entre Node, File et Folder?

\begin{lstlisting}[style=javastyle]
// Create instances of File and Folder
Node node1 = new File("document.txt", 1024);
Node node2 = new Folder("My Documents");

// Store them in a list of type Node
List<Node> nodes = new ArrayList<>();
nodes.add(node1);
nodes.add(node2);

for (Node node : nodes) {
    System.out.println("Name: " + node.name); 
    System.out.println("---");
}

/*
Example Output:
Name: document.txt
---
Name: My Documents
---
*/
\end{lstlisting}

\mcqlist{
 { Node est une interface, et File et Folder sont des classes qui implémentent Node},
 {Node est une classe abstraite, et File et Folder sont des classes qui étendent Node},
 {Node est une classe concrète (non abstraite), et File et Folder sont des classes qui étendent Node},
{
File est une classe, et Node et Folder sont des classes qui étendent File
},
{
Node, File et Folder sont trois classes qui n'ont aucun lien de parenté (héritage).
}
}

\begin{solution}
B, C. Node ne peut pas être une interface car le code accède à un champ d’instance \texttt{node.name}. 
En Java, les interfaces ne peuvent contenir que des attributs \texttt{public static final}, et non des champs dépendant des instances.
\end{solution}

Lesquelles des affirmations suivantes sont correctes:

\mcqlist{
    Une classe Java peut implémenter plusieurs interfaces.,
    Une classe Java peut étendre plusieurs classes.,
     Une classe fille peut accéder aux attributs privés de la classe mère.,
     Une classe déclarée abstract ne peut pas avoir de constructeur. 
}

\begin{solution}
Option A. 
\end{solution}

Soit \texttt{Animal a = new Chien();} (où \texttt{Chien extends Animal}). Si \texttt{Animal} et \texttt{Chien} ont tous deux une méthode \texttt{manger}, laquelle sera appelée par 
\texttt{a.manger()}?

\mcqlist{
Celle de la classe \texttt{Animal},
 Le code ne compile pas car a est de type \texttt{Chien},
Le code ne compile pas car a est de type \texttt{Animal},
Celle de la classe \texttt{Chien}
}

\begin{solution}
Option D.
\end{solution}

\end{Exercice}

\begin{Exercice}[5 minutes]\textbf{Reda}
Soit la spécification de la classe \texttt{HalloweenParty}.

\begin{itemize}
\item Attributs: \begin{itemize} \item Elle possède un attribut \lstinline[style=javastyle]{private Map<String, String> costumes} qui est un dictionnaire associant comme clé le nom d’un invité à une valeur qui est son costume. \end{itemize}
\item Constructeur: \begin{itemize} \item \lstinline[style=javastyle]{public HalloweenParty()}  qui permet de créer une nouvelle instance de la classe avec costumes qui sera initialisé à un dictionnaire vide. \end{itemize}
\item Méthodes: \begin{itemize} \item \lstinline[style=javastyle]{public void addPerson(String name, String costume) }   \newline 
                                        Si le costume n’est pas déjà pris par un invité, l’association (name → costume) est ajoutée (ou mise à jour si name existe déjà), sinon un message de refus est affiché: "Sorry " + name + ", that costume is already taken, you must choose a different one to be added".
                                \item \lstinline[style=javastyle]{private boolean isCostumeTaken(String costume)}  \newline 
                                        Retourne true si costume est une valeur dans le dictionnaire costumes, autrement retourne false.
                                \item \lstinline[style=javastyle]{public String getCostume(String name)}  \newline 
                                        Renvoie le costume associé à name dans le dictionnaire costumes. Si aucune association n’existe, une chaîne de caractères vide est retournée 
\end{itemize}
\end{itemize}

Laquelle de ces affirmations concernant la classe \texttt{HalloweenParty} est vraie?

\mcqlist{
    {On peut accéder à l’attribut \texttt{costumes} depuis une instance de la classe \texttt{HalloweenParty}},
    {La méthode \texttt{getCostume} peut accéder à l’attribut \texttt{costumes}.},
    {
        L'appel \texttt{new HalloweenParty(Map.of("John", "Shrek", "Alicia", "Batman"))} ne lèvera pas d’erreurs
    },
    {
        Après avoir instancié \texttt{HalloweenParty}, si on appelle \texttt{getCostume} avec un nom qui n’est pas une clef de \texttt{costumes}, cela lèvera une erreur
    }
}

\begin{solution}
Option B. 
\end{solution} 

\vspace*{0.5cm}

Laquelle de ces implémentations satisfait la description de la méthode \texttt{addPerson}? \newline \newline 

\begin{lstlisting}[style=javastyle]
// Option A
public void addPerson(String name, String costume) {
    if (!isCostumeTaken(costume)) {
        costumes.put(name, costume);
    } else {
        System.out.println(
            "Sorry " + name
            + ", that costume is already taken, you must choose a different one to be added"
        );
    }
}
\end{lstlisting}

\begin{lstlisting}[style=javastyle]
// Option B
public void addPerson(String name, String costume){
       costumes.put(name, costume);
}
\end{lstlisting}

\begin{lstlisting}[style=javastyle]
// Option C
public boolean addPerson(String name, String costume){
   if(isCostumeTaken(costume)){
         System.out.println("Sorry " + name + ", that costume is already taken, you must choose a different one to be added");
         return;
   } else {
       costumes.put(costume, name);
       return;
  }
}
\end{lstlisting}

\begin{solution}
Option A.
\end{solution}

\vspace*{0.5cm}

Qu’affiche le code suivant? \newline 
\begin{lstlisting}[style=javastyle]
public class Main {
   public static void main(String[] args) {
       HalloweenParty hp = new HalloweenParty();
       hp.addPerson("Richard", "Dinosaur");
       hp.addPerson("Nancy", "Dinosaur");
       hp.addPerson("Richard", "Pilot");
       System.out.println("Richard will be dressed as: " + hp.getCostume("Richard"));
   }
}
\end{lstlisting}

\mcqlist{
    {Cette exécution lèvera une erreur}, 
    {Sorry Nancy, that costume is already taken, you must choose a different one to be added \newline 
    Sorry Richard, that costume is already taken, you must choose a different one to be added \newline 
    Richard will be dressed as: None},
    {
        Sorry Nancy, that costume is already taken, you must choose a different one to be added \newline 
        Richard will be dressed as: Pilot
    },
    {
        Costume already taken, you must choose a different one to be added  \newline 
        Richard will be dressed as: Dinosaur
    }
}

\begin{solution}
Option C.
\end{solution}

\end{Exercice}


\begin{Exercice}[5 minutes]\textbf{Amina}
Qu'affiche le code Java suivant?

\begin{lstlisting}[style=javastyle]
public class Person {
    private String firstName;
    private String lastName;

    public Person(String firstName, String lastName) {
        this.firstName = firstName;
        this.lastName = lastName;
    }

    @Override
    public boolean equals(Object other) {
        if (this == other) return true;
        if (other == null || this.getClass() != other.getClass()) return false;
        Person person = (Person) other;
        return this.firstName.equals(person.firstName)
            && this.lastName.equals(person.lastName);
    }

    // Méthode main pour tester l'identité et l'égalité
    public static void main(String[] args) {
        Person p1 = new Person("Ted", "Mosby");
        Person p2 = new Person("Ted", "Mosby");

        System.out.println((p1 == p2));
        System.out.println(p1.equals(p2));
    }
}
\end{lstlisting}

\mcqlist{
    true \newline 
    true,
    false \newline 
    true,
    true \newline 
    false, 
    false \newline 
    false
}

\begin{solution}
Option B.
\end{solution}

Parmi les affirmations suivantes, lesquelles sont correctes? 

\mcqlist{
    {Une classe abstraite définit un type et fournit une implémentation partielle.},
    {Une interface Java peut contenir du code exécuté, au même titre qu’une classe concrète.},
    {En Java, une classe peut hériter de plusieurs classes mais d’une seule interface.},
    {En Python, les notions de type, classe et interface sont moins strictement différenciées à cause de l’absence de typage statique.}
}
\begin{solution}
Options A, E. 
\end{solution}
\end{Exercice}

\begin{Exercice}[5 minutes]\textbf{Reda}
Soit la classe \texttt{Tournament}

\begin{lstlisting}[style=javastyle]
abstract public class Tournament {
   int nbrOfTeams;

   public Tournament(int  nbrOfTeams) {
       this.nbrOfTeams = nbrOfTeams;
   }

   public abstract String getTournamentType();

   public String getDescription() {
       return "This tournament has " + nbrOfTeams + " teams.";
   }
}
\end{lstlisting}

Quel est l’intérêt de rendre une classe abstraite? (Plusieurs réponses possibles). \newline 

\mcqlist{
{Cela rend son instanciation impossible},
{Permettre de revenir implémenter la classe plus tard},
{Assurer que toutes les classes filles auront la même implémentation des méthodes abstraites},
{Mettre en place une structure commune pour ses classes filles}
}

\begin{solution}
A, D. 
\end{solution}

\vspace*{0.4cm}

Soient les deux classes concrètes qui héritent de \texttt{Tournament}

\begin{lstlisting}[style=javastyle]
public class BasketBallTournament extends Tournament {
   public final int quarterDurationMinutes;
   public BasketBallTournament(int nbrOfTeams, int quarterDurationMinutes) {
       super(nbrOfTeams);
       this.quarterDurationMinutes = quarterDurationMinutes;
   }
   @Override
   public String getTournamentType() {
       return "Basketball tournament. ";
   }

   @Override
   public String getDescription() {
       return super.getDescription() + " Each quarter lasts "  + quarterDurationMinutes + " minutes.";
   }
}
\end{lstlisting}

\begin{lstlisting}[style=javastyle]
public class TennisTournament extends Tournament{
   private final boolean isGrandSlam;
   public TennisTournament(int nbrOfTeams, boolean isGrandSlam) {
       super(nbrOfTeams);
       this.isGrandSlam = isGrandSlam;
   }
   @Override
   public String getTournamentType() {
       return "Tennis tournament. ";
   }
   @Override
   public String getDescription() {
       String additionalInfo;
       if (isGrandSlam) {
           additionalInfo = " This tournament is a Grand Slam.";
       } else {
           additionalInfo = " This tournament is not a Grand Slam.";
       }
       return super.getDescription() + additionalInfo;
   }
}
\end{lstlisting}

Aurions-nous eu une erreur si on retirait l’implémentation de \texttt{getDescription} dans \texttt{TennisTournament}?

\mcqlist{
    {Oui, car la classe mère est abstraite, et il faut implémenter chaque méthode},
    {Non, car nous avons déjà une implémentation dans la classe mère},
    {Oui, car le type de retour de la fonction est String},
    {Non, car on peut implémenter une classe comme on le souhaite}
}

\begin{solution}
Option B.
\end{solution}

Aurions-nous eu une erreur si on retirait l’implémentation de \texttt{getTournamentType} dans \texttt{BasketballTournament}?

\mcqlist{
    {Oui, car la classe mère est abstraite, et il faut implémenter chaque méthode},
    {Non, car le @Override indique qu’elle peut être retirer},
    {Oui, car la méthode est abstraite dans la classe mère},
    {Non, car on peut implémenter une classe comme on le souhaite}
}

\begin{solution}
Option C. 
\end{solution}


Qu’affiche le programme suivant? \newline

\begin{lstlisting}[style=javastyle]
public class Main {
   public static void main(String[] args) {
       Tournament tournament1 = new BasketBallTournament( 32, 10);
       Tournament tournament2 = new TennisTournament( 128, true);

       System.out.println("Tournament 1: " + tournament1.getTournamentType() + tournament1.getDescription());
       System.out.println("Tournament 2: " + tournament2.getTournamentType() + tournament2.getDescription());
   }
}
\end{lstlisting}
\mcqlist{
    {Ce programme va lever une erreur},
    {Tournament 1: This tournament has 32 teams. \newline 
    Tournament 2: This tournament has 128 teams. },
    {Tournament 1: Basketball tournament. This tournament has 32 teams. Each quarter lasts 10 minutes. \newline 
    Tournament 2: Tennis tournament. This tournament has 128 teams. This tournament is a Grand Slam.},
    {
Tournament 1: Basketball tournament. This tournament has 32 teams. \newline 
Tournament 2: Tennis tournament. This tournament has 128 teams. 
    }
}

\begin{solution}
Option C. 
\end{solution}

\vspace*{0.3cm}

Comment peut-on accéder à l’attribut \texttt{quarterDurationMinutes} de \texttt{tournament1}?

\mcqlist{
    {Il faut faire un type casting},
    {Il n’est pas possible d’accéder à cet attribut},
    {Il faut implémenter un getter dans la classe \texttt{Tournament}},
    {Il faut utiliser \texttt{tournament1.quarterDurationMinutes}}
}

\begin{solution}
Option A.
\end{solution}

\end{Exercice}

